\documentclass{article}
\usepackage{amsmath}
\begin{document}
\section*{The same conversion budget can favor different ADC schedules}
The ADS1015 comparison fixed the available row menu. This extension
uses that exact menu but models the raw pin levels as moving between
conversions. It is a finite acquisition model using specified time
slots and qualified stable signs, not a hardware timing estimate.
The target is the normalized balanced direction
$g=(1,1,-1,-1)$ at a decision cut $T$. This target needs three
conversions even on frozen data.

\paragraph{Model and sufficient data.}
A conversion of an allowed singleton or differential row $q_j$
at time $t_j$ returns
\[
    y_j=q_jv(t_j)+b(t_j)\sum_i q_{ji}+e_j .
\]
Each raw pin trajectory $v_i$ is independently $L_i$-Lipschitz.
The common reference error is $D$-Lipschitz and has amplitude at
most $B$; independent conversion errors satisfy $|e_j|\le\eta_j$.
All these bounds are simultaneous hard supports with no additional
coupling. Absolute levels of the four pins are free in this linear
model. Quantizer clipping, physical rail and common-mode bounds,
settling, conversion aperture, unstable sign interpretation and
calibration uncertainty outside these terms are excluded. The
budget allows at most three observations at the fixed slots
$t_0,t_1,t_2$, with $T\ge t_2$.

A finite uniform error requires the chosen rows to span $g$:
otherwise a constant shift in their nullspace changes $gv(T)$
without changing any data. The target has no two-row representation,
so any finite solution here uses three independent rows with the
unique coefficients $\alpha$ given by $g=\sum_j\alpha_jq_j$.
Repeats, dependencies, and stopping short yield no finite bound at
this target under the free-level assumption; they are not pruned
merely to simplify the enumeration.

\paragraph{Exact radius, not just a propagated error bound.}
Use $\sum_j\alpha_j y_j$ as the center. Put $t_3=T$.
For pin $i$, the coefficient of its $j$th interval increment
has magnitude $|\sum_{k\le j}\alpha_kq_{ki}|$ by telescoping. The
exact support of the difference between the center and $gv(T)$ is
the sum of
\begin{align*}
 R_{\rm pin}&=\sum_i L_i\sum_{j=0}^2
    \left|\sum_{k\le j}\alpha_kq_{ki}\right|(t_{j+1}-t_j),\\
 R_{\rm noise}&=\sum_j|\alpha_j|\eta_j
\end{align*}
and the support of the bias weights
$\beta_j=\alpha_j\sum_i q_{ji}$. The bias sum is zero. Valid
three-row representations have either zero singletons or two:
one singleton would have coefficient zero, which would imply a
two-row representation, and three singletons cannot touch all
four nonzero coordinates. In the all-differential case bias
support is zero. With two singleton slots $u,v$, it is exactly
$|\beta_u|\min(2B,D|t_u-t_v|)$. Each term can attain its
worst endpoint jointly by choosing independent extremal pin
increments, a centered bias pair linked by a linear admissible
path, and extremal read errors.

For any fixed observation vector, independent rows let a constant
shift of the pin trajectories produce that vector without changing
the error; solve the full-row-rank linear system for the shift.
Thus each extremal error is achievable at the same transcript.
The radius is minimax against arbitrary estimators, not just
linear error propagation. This also compares deterministic adaptive
row choice on the same fixed time budget: follow an all-zero
output transcript to identify the chosen rows and then use the
two extremal trajectories, shifted to produce exactly that
transcript. The controller would choose those same rows in both
histories. An invalid row set instead has a constant nullspace
obstruction. A fixed minimizing design attains the lower bound.
This argument does not evaluate adaptive acquisition times,
randomized risk objectives, clipped domains or correlated signs.

\paragraph{A sharp reversal at equal conversion count.}
At times $(0,1,2)$ and cut $T=2$, take $L=(1,1,1,4)$
and uniform independent conversion error $\eta$.
The best mixed path reads singleton 0, then singleton 2,
then differential $1-3$ (signs $+,-,+$); its pin radius is $3$,
with the two singleton bias terms one slot apart. The best
all-differential path reads $0-3$, then $2-3$, then $1-3$
(signs $+,-,+$); its pin radius is $7$, because an older
occurrence of the moving pin 3 matters before cancellation.
The following minima of $R_{\rm pin}$ and the noise
multipliers enumerate the five unordered spanning menus:
\[
\begin{array}{c|cc}
\text{rows}& \min R_{\rm pin} & \sum |\alpha_j|\\ \hline
 01,03,23 &9&4\\
 01,13,23 &9&4\\
 03,13,23 &7&3\\
 03,S_1,S_2 &3&3\\
 13,S_0,S_2 &3&3
\end{array}
\]
Here each two-digit label is an ordered differential and
$S_i$ denotes a singleton. The primary 30 finite schedules include
all permutations of these five menus. In each mixed minimum an
adjacent singleton pair is attainable, the shortest spacing. The
resulting optimum is
\[
  3\eta+\min\{3+\min(2B,D),\,7\}.
\]
The mixed path wins below effective drift $4$, the differential
path wins above it, and they tie at equality. Amplitude can keep
the mixed path preferable even if reference variation is fast.

The replay checks all 585 ordered mode sequences of length at most
three including repetitions (and zero observations), then 810
independent primal set-membership LPs on the exact native MUX menu.
The LP computes indistinguishable target extrema directly from
a zero transcript; it is not given coefficient or error formulas.
It agrees. This resolves a choice at the modeled cut, but does not
instantiate real conversion windows, prove a particular motor
force, or use the independent non-radio evaluation.
\end{document}
